\begin{conclusion}[例外幂和的三项主导展开与显式余项界]\label{con:xi-terminal-replica-softcore-exceptional-moments-exponential-hierarchy}
取整数 \(m\ge 3\)，并定义
\[
A_{m,0}:=2^{m},\qquad
A_{m,1}:=3\cdot 2^{m-2},\qquad
A_{m,2}:=5\cdot 2^{m-3}.
\]
则对一切 \(q\ge 2\) 存在严格分解
\[
\boxed{\;
S_m(q)=2^{-m}\Big(A_{m,0}^{q}+mA_{m,1}^{q}+mA_{m,2}^{q}+\Omega_m(q)+R_m(q)\Big)\;}
\]
其中余项 \(R_m(q)\) 满足显式界
\[
0\le R_m(q)\le (2^{m}-2m-2)\,\bigl(9\cdot 2^{m-4}\bigr)^{q}.
\]
从而得到相对误差展开
\[
\frac{S_m(q)}{2^{m(q-1)}}=
1+m\left(\frac34\right)^{q}
+m\left(\frac58\right)^{q}
+O_m\!\left(\left(\frac{9}{16}\right)^{q}\right),
\qquad (q\to\infty),
\]
其中 \(O_m\) 的隐常数仅依赖于 \(m\)。
\end{conclusion}
\begin{proof}
由结论 \ref{con:xi-terminal-replica-softcore-exceptional-power-sum-word-trace-sympower}，
\[
2^mS_m(q)=\Omega_m(q)+\sum_{\substack{w\in\{J,K\}^m\\ w\neq K^m}}\bigl(\Tr M(w)\bigr)^q,
\qquad M(w)=A_1\cdots A_m.
\]
将求和按字母 \(K\) 的出现次数 \(r=\#\{t:A_t=K\}\) 分层。
\par
\textbf{(i) \(r=0\).}
仅有字词 \(w=J^m\)，且 \(\Tr(J^m)=2^m=A_{m,0}\)。
\par
\textbf{(ii) \(r=1\).}
恰有 \(m\) 个字词，为 \(KJ^{m-1}\) 的循环移位；由结论 \ref{con:xi-terminal-replica-softcore-fixed-zero-count-trace-extremal}（取 \(r=1\)）知该层词迹恒为
\(\Tr M(w)=3\cdot 2^{m-2}=A_{m,1}\)。
\par
\textbf{(iii) \(r=2\).}
由注 \ref{rem:xi-terminal-replica-softcore-trace-moment-cycle-subgraph-word-identification}，
\(\Tr M(w)\) 等于相应循环子图 \(C_m[S]\) 的独立集计数 \(I(C_m[S])\)。
当 \(S\) 为两条相邻边时 \(I(C_m[S])=5\cdot 2^{m-3}=A_{m,2}\)，且这类 \(S\) 恰有 \(m\) 个；
因此相邻型字词总贡献为 \(mA_{m,2}^{q}\)。
\par
\textbf{(iv) 余项界。}
对其余全部字词（仍满足 \(w\neq K^m\) 且不属于 (i)--(iii)），若 \(r\ge 3\)，由
\(
\Tr M(w)\le F_{r+3}\,2^{m-r-1}\le 2^{r}\,2^{m-r-1}=2^{m-1}<9\cdot 2^{m-4}
\)
得 \(\Tr M(w)<9\cdot 2^{m-4}\)；
而 \(r=2\) 的非相邻型给出 \(I(C_m[S])=9\cdot 2^{m-4}\)，故亦有 \(\Tr M(w)\le 9\cdot 2^{m-4}\)。
令 \(R_m(q)\) 为上述其余字词之 \((\Tr M(w))^q\) 的总和。
该类字词个数为 \((2^m-1)-(1+m+m)=2^m-2m-2\)，因此
\[
0\le R_m(q)\le (2^{m}-2m-2)\,\bigl(9\cdot 2^{m-4}\bigr)^{q}.
\]
\par
合并各层贡献并除以 \(2^m\) 即得严格分解。
最后，将分解式两侧除以 \(2^{m(q-1)}=2^{-m}A_{m,0}^q\) 得
\[
\frac{S_m(q)}{2^{m(q-1)}}=
1+m\left(\frac{A_{m,1}}{A_{m,0}}\right)^q
+m\left(\frac{A_{m,2}}{A_{m,0}}\right)^q
+\frac{\Omega_m(q)}{2^{mq}}
+\frac{R_m(q)}{2^{mq}}.
\]
其中 \(\frac{A_{m,1}}{A_{m,0}}=\frac34\) 且 \(\frac{A_{m,2}}{A_{m,0}}=\frac58\)。
另一方面，\(\Omega_m(q)\) 的闭式（结论 \ref{con:xi-terminal-replica-softcore-lucas-pell-omega-unified}）给出
\(|\Omega_m(q)|\le C_m\,\varphi^{mq}\)（常数 \(C_m>0\) 仅依赖于 \(m\)），从而
\[
\frac{|\Omega_m(q)|}{2^{mq}}\le C_m\left(\frac{\varphi}{2}\right)^{mq}
\le C_m\left(\frac{9}{16}\right)^q
\qquad(m\ge 3).
\]
结合 \(\frac{R_m(q)}{2^{mq}}\le (2^m-2m-2)\left(\frac{9}{16}\right)^q\)，即得所示相对误差展开。证毕。
\end{proof}

\begin{conclusion}[全迹矩的三项主导展开与显式余项界]\label{con:xi-terminal-replica-softcore-trace-moments-exponential-hierarchy}
保持结论 \ref{con:xi-terminal-replica-softcore-exceptional-moments-exponential-hierarchy} 的记号（取 \(m\ge 3\) 并沿用 \(A_{m,0},A_{m,1},A_{m,2}\) 与余项 \(R_m(q)\)）。
则对一切 \(q\ge 2\) 成立严格分解
\[
\boxed{\;
\Tr\!\bigl((T^{(q)})^m\bigr)
=2^{-m}\Big(A_{m,0}^{q}+mA_{m,1}^{q}+mA_{m,2}^{q}+L_m^{q}+R_m(q)\Big)\;}
\]
并且余项界
\[
0\le R_m(q)\le (2^{m}-2m-2)\,\bigl(9\cdot 2^{m-4}\bigr)^{q}
\]
仍然成立。
因此
\[
\frac{\Tr((T^{(q)})^m)}{2^{m(q-1)}}
=1+m\left(\frac34\right)^{q}
+m\left(\frac58\right)^{q}
+O_m\!\left(\left(\frac{9}{16}\right)^{q}\right),
\qquad (q\to\infty),
\]
其中 \(O_m\) 的隐常数仅依赖于 \(m\)。
\end{conclusion}
\begin{proof}[证明要点]
由结论 \ref{con:xi-terminal-replica-softcore-trace-exceptional-difference}，
\(
\Tr((T^{(q)})^m)=S_m(q)+2^{-m}(L_m^q-\Omega_m(q))
\)；
将结论 \ref{con:xi-terminal-replica-softcore-exceptional-moments-exponential-hierarchy} 的
\(
S_m(q)=2^{-m}\bigl(A_{m,0}^{q}+mA_{m,1}^{q}+mA_{m,2}^{q}+\Omega_m(q)+R_m(q)\bigr)
\)
代入即可抵消 \(\Omega_m(q)\)，得到所示分解与同一余项界。相对误差展开由两侧同除以 \(2^{m(q-1)}\)，再用 \(A_{m,1}/A_{m,0}=3/4\)、\(A_{m,2}/A_{m,0}=5/8\) 以及 \((L_m/2^m)^q\le (\varphi/2)^{mq}\le (9/16)^q\)（\(m\ge 3\)）并结合 \(R_m(q)\) 的界给出。证毕。
\end{proof}

\endinput

